% Figure: thermal and velocity boundary layers growing over a flat plate.
% Free stream U_inf, T_inf approaches the leading edge; the boundary layer
% thickness grows like sqrt(x) in the laminar region, then thickens and
% transitions to turbulent. Plate held at T_s.
\begin{figure}[htbp]
\centering
\begin{tikzpicture}[scale=1.0, >=latex]
% plate
\fill[enggray!25] (0,-0.35) rectangle (10,0);
\draw[thick] (0,0) -- (10,0);
\node[font=\scriptsize] at (5,-0.62) {heated plate at $T_s$};
% free-stream arrows
\foreach \y in {1.4,1.9,2.4}{
  \draw[->, engblue, thick] (-1.1,\y) -- (-0.2,\y);
}
\node[font=\scriptsize, engblue, anchor=east] at (-1.15,1.9) {$U_\infty,\,T_\infty$};
% laminar boundary-layer edge: delta ~ sqrt(x), from 0 to x=5
\draw[engred, thick, samples=60, domain=0.02:5, variable=\x]
  plot ({\x}, {0.62*sqrt(\x)});
% transition marker
\draw[enggray, dashed] (5,0) -- (5,2.2);
\node[font=\scriptsize, enggray] at (5,2.45) {transition};
% turbulent region: thicker, grows faster ~ x^{0.8}, continue to x=10
\draw[engred, thick, samples=60, domain=5:10, variable=\x]
  plot ({\x}, {1.275*0.30*(\x^0.8)});
\node[font=\scriptsize, engred, anchor=west] at (9.0,2.7) {$\delta(x)$};
% velocity profile glyphs at two stations
\foreach \xs/\dl in {2/0.88,8/2.0}{
  \draw[engblue, thin] (\xs,0) -- (\xs,\dl);
  \foreach \f in {0.25,0.5,0.75,1.0}{
    \draw[->, engblue, thin] (\xs,{\f*\dl}) -- ({\xs + 0.7*sqrt(\f)},{\f*\dl});
  }
}
\node[font=\scriptsize, engblue] at (2.0,1.15) {laminar};
\node[font=\scriptsize, engblue] at (8.0,2.3) {turbulent};
% x axis
\draw[->] (0,-0.95) -- (10.2,-0.95) node[right, font=\scriptsize] {$x$};
\node[font=\scriptsize] at (0,-1.2) {leading edge};
\end{tikzpicture}
\caption[Boundary layer growth over a heated flat plate]{The boundary
layer growing over a heated flat plate. The free stream arrives at
velocity $U_\infty$ and temperature $T_\infty$; the layer of fluid
slowed and warmed by the plate thickens with distance, as
$\delta \sim \sqrt{x}$ in the laminar region near the leading edge. Past
a transition Reynolds number near $5\times 10^5$ the layer becomes
turbulent, thickens faster, and mixes more vigorously, which raises the
local convection coefficient. The velocity profiles sketched at two
stations show the no-slip condition at the wall and the approach to the
free stream at the boundary-layer edge. The local Nusselt-number
correlations of section 5.3 follow this same $\sqrt{x}$ scaling in the
laminar region and the steeper $x^{4/5}$ scaling once turbulent.}
\label{fig:vol04ch05:boundary-layer}
\end{figure}
